% Quelle: 1. Übung Lösungen.pdf
% Art: Übung mit Lösungen
% Themen: Teilbarkeitsregeln, Teilbarkeit durch 9, 11 und 7, Rundenturniere, Rundenplan, Blockpläne, BIBD
\section{Übung 1 (mit Lösungen)}
\subsection{Aufgaben}

\subsubsection*{Teilbarkeitsregeln}

\begin{aufgabe}[1]
Sei $x = 299\,792$. \quad (Geschwindigkeit des Lichtes in [km / s])\\
Ist $x$ durch 9 oder durch 11 teilbar?
\end{aufgabe}

\begin{aufgabe}[2]
Entwickeln Sie eine Regel für die Division durch 7. Beachten Sie dabei
\[ 1001 = 7 \cdot 11 \cdot 13 \]
Entscheiden Sie mit Hilfe dieser Regel, ob\\
$x = 299\,792$ oder $x = 1000\;0000\;0000\;0000\;0001$ durch 7 teilbar sind.
\end{aufgabe}

\subsubsection*{Rundenturniere}

\begin{aufgabe}[3]
Stellen Sie einen Rundenplan für $2m = 8$ Mannschaften an 7 Tagen auf.
\end{aufgabe}

\begin{aufgabe}[4]
Gegen wen spielt Mannschaft $2m$ an Tag $r$?
\end{aufgabe}

\subsubsection*{Blockpläne, BIBD}

\begin{aufgabe}[5]
Beweisen Sie die Gleichungen
\[ b\,k = v\,r, \qquad\qquad r\,(k-1) = \lambda\,(v-1) \]
für BIBDs.\\
Hinweis: Für die zweite Gleichung lohnt es sich die 2-elementigen Teilmengen im BIBD zu zählen, und auf das Resultat die erste Gleichung anzuwenden.
\end{aufgabe}

\subsection{Ü1 Lösungen}

\subsubsection*{Teilbarkeitsregeln}

\begin{loesung}
\textbf{1\_1}
\begin{align*}
x &= 299\,792\\
x &\equiv x_n + \dots + x_1 + x_0 \pmod{9}\\
  &= 2 + 9 + 9 + 7 + 9 + 2\\
  &\equiv 2 + 7 + 2 \equiv 2 \pmod{9}\\
\Rightarrow\quad & x = 299\,792 \text{ ist nicht durch 9 teilbar}
\end{align*}
\begin{align*}
x &\equiv (-1)^5\, x_5 + \dots - x_1 + x_0 \pmod{11}\\
  &= -2 + 9 - 9 + 7 - 9 + 2\\
% KORRIGIERT: "= -2 (mod 11)" statt "\equiv"
  &\equiv -2 \not\equiv 0 \pmod{11}\\
\Rightarrow\quad & x = 299\,792 \text{ ist nicht durch 11 teilbar}
\end{align*}
\end{loesung}

\begin{loesung}
\textbf{1\_2}
\begin{align*}
% KORRIGIERT: Malpunkt fehlte bei y_1 1000; Begründung 1000 = 1001 - 1 ergänzt
x &= y_m \cdot 1000^m + \dots + y_1 \cdot 1000 + y_0\\
  &\equiv (-1)^m\, y_m + \dots - y_1 + y_0 \pmod{7} && \text{(da } 1000 = 1001 - 1 \equiv -1 \pmod{7}\text{)}\\[1ex]
x &= 299\,792\\
% KORRIGIERT: "(mod 7)" fehlte, Rest 3 nicht als Kongruenz notiert
  &\equiv -299 + 792 = 493 = 7 \cdot 70 + 3 \equiv 3 \pmod{7}\\
\Rightarrow\quad & x = 299\,792 \text{ ist nicht durch 7 teilbar}\\[1ex]
x &= 1000\;0000\;0000\;0000\;0001 = 10^{19} + 1 = 10 \cdot (10^3)^6 + 1\\
  &\equiv 10 \cdot (-1)^6 + 1 \pmod{7}\\
% KORRIGIERT: Rest 11 nicht reduziert
  &\equiv 11 \equiv 4 \pmod{7}\\
\Rightarrow\quad & x = 1000\;0000\;0000\;0000\;0001 \text{ ist nicht durch 7 teilbar.}
\end{align*}
\end{loesung}

\subsubsection*{Rundenturniere}

\begin{loesung}
\textbf{1\_3} Stellen Sie einen Rundenplan für $2m = 8$ Mannschaften an 7 Runden auf.
\begin{align*}
% KORRIGIERT: "x \neq y" und "(mod 2m-1)" in der zweiten Zeile fehlten
& x + y \equiv r \pmod{2m-1} \quad \text{für } 1 \le x, y \le 2m-1,\ x \neq y,\\
& \text{falls } x + x \equiv r \pmod{2m-1}\text{, spielt } x \text{ gegen } 2m.
\end{align*}
\begin{center}
\begin{tabular}{ccccccc}
\toprule
$r=1$ & $r=2$ & $r=3$ & $r=4$ & $r=5$ & $r=6$ & $r=7$\\
\midrule
1\ 7 & 1\ 8 & 1\ 2 & 1\ 3 & 1\ 4 & 1\ 5 & 1\ 6\\
2\ 6 & 2\ 7 & 3\ 7 & 2\ 8 & 2\ 3 & 2\ 4 & 2\ 5\\
3\ 5 & 3\ 6 & 4\ 6 & 4\ 7 & 5\ 7 & 3\ 8 & 3\ 4\\
4\ 8 & 4\ 5 & 5\ 8 & 5\ 6 & 6\ 8 & 6\ 7 & 7\ 8\\
\bottomrule
\end{tabular}
\end{center}
\end{loesung}

\begin{loesung}
\textbf{1\_4} Gegen wen spielt Mannschaft $2m = 8$ in Runde $r$?
\begin{align*}
x + x = 2 \cdot x &\equiv r && \pmod{2m-1}\\
2m &\equiv 1 && \pmod{2m-1}\\
% KORRIGIERT: "x = r m (mod 2m-1)" statt "\equiv"; Begründung für die Inverse von 2 ergänzt
\Rightarrow\quad 2^{-1} &\equiv m && \pmod{2m-1} && \text{(existiert, da } \ggT(2, 2m-1) = 1\text{)}\\
\Rightarrow\quad x &\equiv r \cdot m && \pmod{2m-1} && \text{(beide Seiten mal } 2^{-1}\text{)}
\end{align*}
z.\,B.
\begin{center}
\begin{tabular}{lll}
$r = 1$ & $x = 1 \cdot 4 = 4$ & \\
$r = 2$ & $x = 2 \cdot 4 = 8$ & $\equiv 1 \pmod{7}$\\
$r = 3$ & $x = 3 \cdot 4 = 12$ & $\equiv 5 \pmod{7}$\\
$r = 4$ & $x = 4 \cdot 4 = 16$ & $\equiv 2 \pmod{7}$\\
$r = 5$ & $x = 5 \cdot 4 = 20$ & $\equiv 6 \pmod{7}$\\
$r = 6$ & $x = 6 \cdot 4 = 24$ & $\equiv 3 \pmod{7}$\\
% KORRIGIERT: "\equiv 7 (mod 7)" ohne Hinweis auf 28 \equiv 0 und Vertreter 7
$r = 7$ & $x = 7 \cdot 4 = 28$ & $\equiv 0 \equiv 7 \pmod{7}$ (Vertreter aus $1, \dots, 7$)
\end{tabular}
\end{center}
\end{loesung}

\subsubsection*{Blockpläne, BIBD}

\begin{loesung}
\textbf{1\_5} Beweisen Sie die Gleichungen
\[ b\,k = v\,r, \]
für BIBDs.\\
Zählen die Varietäten im BIBD: \quad $b\,k = v\,r$.
\[ r\,(k-1) = \lambda\,(v-1) \]

\textbf{1. Lösung:}\\
% KORRIGIERT: Malpunkte ergänzt; Begründungen für Einsetzen und Kürzen durch v fehlten
Zählen die 2-elementigen Teilmengen im BIBD: \quad $b \cdot \binom{k}{2} = \lambda \cdot \binom{v}{2}$.
\begin{align*}
\Rightarrow\quad & \lambda\, v\,(v-1) = b\,k\,(k-1) = v\,r\,(k-1) && \text{(mal 2; dann } b\,k = v\,r\text{)}\\
\Rightarrow\quad & \lambda\,(v-1) = r\,(k-1) && \text{(geteilt durch } v > 0\text{)}
\end{align*}
\textbf{2. Lösung:} Zählen in wie vielen Paaren ein $a \in V$ im BIBD vorkommt.
\[ r\,(k-1) = \lambda\,(v-1) \]
\end{loesung}
